\chapter*{Abstract}
\addcontentsline{toc}{chapter}{Abstract}

Cells cultured in organ-on-chip devices respond to the fluid shear stress on their apical surface, and they need oxygen. Delon \emph{et al.} (2020) proposed a Hele-Shaw chip whose width grows hyperbolically, $w(x)=w_1L/(L-x)$, so that the wall shear stress falls linearly along the channel. A single device then exposes cells to a whole range of shear stresses. This project reproduces that device numerically and extends it to oxygen transport.

The fluid domain has a \SI{1}{mm} inlet, a hyperbolic section ($L=\SI{50}{mm}$) that widens to \SI{10}{mm} at $x=\SI{45}{mm}$, a \ang{45} converging outlet and a uniform height of \SI{145}{\micro\metre}. Steady, laminar, three-dimensional ANSYS Fluent simulations were run at 0.22, 1 and \SI{120}{\uLmin}. They were compared with three levels of theory derived here: the parallel-plate lubrication solution used by the reference paper (T1), its exact rectangular-duct correction for side-wall drag (T2), and a two-dimensional depth-averaged Hele-Shaw model of the complete planform (T3).

The flow is in the Stokes regime ($\Rey_{D_h}\le3.65$). The pressure drop is \SI{1.657}{Pa} per \ulmin\ and every interior quantity scales with $Q$ to within \SI{0.08}{\percent}. At the four stations S1--S4 the centre-line wall shear agrees with the published formula (T1) to within $\pm\SI{3.3}{\percent}$, and the shear range reproduces the values reported by Delon \emph{et al.}: \SIrange{0.0102}{0.0014}{\dyncm} against $\sim$\SIrange{0.01}{0.001}{\dyncm} at \SI{0.22}{\uLmin}. This agreement is partly a cancellation, however. The CFD correctly contains the side-wall drag that T1 omits (\SI{+9}{\percent} at S1), and it is uniformly \SI{5.2(4)}{\percent} below the side-wall-corrected theory. The most probable cause is too few cells across the gap; a mesh-independence study has not yet been performed. Ratios that cancel this offset agree with T3 to within \SI{1}{\percent}. At S5 the shear is \SI{46.5}{\percent} above T1 because S5 lies on the outlet corner, which T3 predicts (\SI{+54}{\percent}). Within $\pm\SI{5}{\percent}$, the shear is uniform over \SI{70}{\percent} (S1) to \SI{91}{\percent} (S4) of the width, bounded by side-wall layers \SIrange{0.15}{0.17}{mm} thick (theory \SI{0.129}{mm}). This uniformity does not depend on the flow rate.

For oxygen, a convection--diffusion model with Michaelis--Menten uptake at the cell-covered floor was derived. It was reduced to a cup-mixing balance in which the geometry enters only through the upstream cell area $A(x)=w_1L\ln[L/(L-x)]$, and solved in closed form for impermeable, PDMS and saturated walls. With illustrative uptake parameters and impermeable walls, convection alone cannot feed the cells below $Q_{\min}=\SI{1.55}{\uLmin}$. The model predicts anoxia beyond $x=\SI{13.9}{mm}$ at \SI{0.22}{\uLmin} and beyond \SI{38.7}{mm} at \SI{1}{\uLmin}, which is exactly the low-shear region of the chip. A \SI{5}{mm} PDMS lid can supply about 25 times the maximum uptake and keeps the oxygen concentration at the cells at or above \SI{95}{\percent} of saturation. The transverse oxygen drop across the gap is below \SI{1.1}{\percent}. Shear and oxygen availability are therefore confounded at low flow rates unless the walls are gas-permeable. These oxygen results are model predictions; they still have to be tested by oxygen CFD and by measurements.

\vspace{3mm}
\noindent\textbf{Keywords:} microfluidics; Hele-Shaw flow; wall shear stress; organ-on-chip; lubrication theory; CFD validation; oxygen transport; Michaelis--Menten kinetics.
